\subsection{DC motor 1\label{AufgGM1}}
\Aufgabe{
    An externally excited DC machine has the following information on its nameplate:
    \begin{itemize}
        \item Rated power: $41,8\,\text{kW}$
        \item Rated speed: $1900\,\frac{1}{\text{min}}$
        \item Anchor voltage: $440\,\text{V}$
        \item Anchor current: $100\,\text{A}$
        \item Excitement: $240\,\text{V}$
        \item Excitation current: $10\,\text{A}$
        \item Idling speed: $2000\,\frac{1}{\text{min}}$
    \end{itemize}

Note: The rated power on a motor nameplate always refers to the mechanical power output at the rated point.

    \begin{itemize}
        \item[a)]
        What is the torque of the machine at the rated point?
        \item[b)]
        Calculate the efficiency of the machine at the rated point.
        \item[c)]
        Calculate the product of the excitation flux and the armature constant $K\cdot \Phi$
		\item[d)]
        How large is the anchor resistance $R_A$?
	    \end{itemize}
}

\Loesung{
	\begin{itemize}

		\item[a)]
        The rated torque is directly derived from the rated power and rated speed:
		\begin{align*}
		    P_\text{mech} &= M\cdot \omega\\
		    M &= \frac{P_\text{mech}}{\omega} = \frac{41,8\,\text{kW}}{2\pi\cdot 1900\,\frac{1}{\text{min}} \cdot \frac{1\,\text{min}}{60\,\text{s}}} = 210,08\,\text{Nm}
		\end{align*}
		\item[b)]
        In motor operation, the efficiency is the quotient of mechanical power relative to electrical power (see equation \ref{GlLeistung}). Both the armature power and the excitation power must be taken into account here.
		\begin{eqa}
			\eta &= \frac{P_\text{mech}}{P_\text{elektr}} = \frac{41,8\,\text{kW}}{440\,\text{V}\cdot 100\,\text{A} + 240\,\text{V}\cdot 10\,\text{A}} = 90,08\,\%\nonumber
		\end{eqa}

		\item[c)]
        The induced voltage is calculated using equation \ref{GlInduzierteSpannungGM}. When idling, the induced voltage must correspond to the armature voltage. The same can be seen from equation \ref{GlfremderregteGM1}, since the torque must also be zero when idling.
		\begin{eqa}
			U_q &= K \cdot \Phi \cdot \omega\tag{\ref{GlInduzierteSpannungGM}}\\
			K\cdot \Phi &= \frac{U_q}{\omega} = \frac{440\,\text{V}}{2\pi\cdot 2000\,\frac{1}{\text{min}} \cdot \frac{1\,\text{min}}{60\,\text{s}}} = 2,1\,\text{Vs}\nonumber
		\end{eqa}
		\item[d)]
        There are two possible approaches here. Equation \ref{GlfremderregteGM1} can be transformed to the anchor resistance:
		\begin{eqa}
			n &= \frac{U_A}{2\pi K\cdot \Phi} - \frac{R_A\cdot M_i}{2\pi(K\cdot \Phi)^2}\tag{\ref{GlfremderregteGM1}}\\
			R_A &= \left(\frac{U_A}{2\pi K\cdot \Phi} - n\right) \cdot \frac{2\pi(K\cdot \Phi)^2}{M_i}\nonumber\\
			&= \left(\frac{440\,\text{V}}{2\pi\cdot 2,1\,\text{Vs}} - \frac{1900}{60\,\text{s}}\right)\cdot \frac{2\pi\cdot (2,1\,\text{Vs})^2}{210,08\,\text{Nm}} = 220\,\text{m}\Omega\nonumber
		\end{eqa}
        The other calculation variant uses the equivalent circuit diagram in Figure \ref{AbbGMFremderregt}. In nominal operation, we can calculate the source voltage $U_q$. This must be smaller than the source voltage in no-load operation.
		\begin{eqa}
			U_q &= K \cdot \Phi \cdot \omega \tag{\ref{GlInduzierteSpannungGM}}\\
			U_{q,n}&=  2,1\,\text{Vs} \cdot 2\pi\cdot 1900\,\frac{1}{\text{min}} \cdot \frac{1\,\text{min}}{60\,\text{s}} = 418\,\text{V}\nonumber
		\end{eqa}
        The voltage drop across resistor $R_A$ must be equal to the difference between the armature voltage and the induced voltage through the mesh circulation. Since the armature current is known in nominal operation, the resistance can be calculated using Ohm's law.
		\begin{equation*}
		    R_A = \frac{U_A - U_q}{I_A} = \frac{440\,\text{V} - 418\,\text{V}}{100\,\text{A}} = 220\,\text{m}\Omega
		\end{equation*}
	\end{itemize}
}